\begin{table}[h]
\centering
\caption{Determinants of Highway Placement - PPML on Discretionary Sample}
\label{tab:ppml_results}
\begin{threeparttable}
\begin{tabular*}{0.7\textwidth}{@{\extracolsep{\fill}}l*{1}{r}}
\toprule
 & Coefficient \\
\midrule
Residential & \makecell[tr]{-0.744 \\ (1.159)} \\
Black & \makecell[tr]{1.123 \\ (1.186)} \\
Residential x Black & \makecell[tr]{1.282 \\ (7.446)} \\
Pseudo R-squared & 0.435 \\
Observations & 5301 \\
\bottomrule
\end{tabular*}
\begin{tablenotes}[flushleft]
\footnotesize
\item This table contains estimates of the impact of residential zoning and majority-Black status on the likelihood of highway placement, estimated using a Poisson Pseudo-Log-Linear model. The sample is restricted to a subset of grid squares designated as discretionary. The discretionary sample excludes grid squares that are intersected by a highway in 1940 or are directly adjacent to a highway in 1940. Standard errors are reported in parenthesis and estimated using a bootstrap procedure with 500 draws. The model includes controls for housing, geographic, and demographic characteristics, as well as city fixed effects. This model also includes a measure of a square's geographic suitability for highway placement, as estimated by a CNN model. * p<0.10, ** p<0.05, *** p<0.01
\end{tablenotes}
\end{threeparttable}
\end{table}
